\documentclass[../main.tex]{subfiles}
\graphicspath{{\subfix{../images/}}}
\begin{document}
%\section{Shells and Stuff}
\subsection{What are shells and membranes? The Definitions}
Most structural objects/elements can be reduced to a physical line with a certain cross-section. The description of stresses in these objects can be done by describing
 “stress resultants” - the normal force, transverse shearing forces, bending moments, and the torque - in the cross section. In contrast to these objects,
 certain structural elements are made to enclose some space making it impossible for them to be reduced to a line. They are better described as a curved surface or a plane.
 These elements can be broadly broken down in two categories - plates and shells [1].

A shell is a curved, thin-walled structure with the thickness being significantly smaller than the other dimensions [2]. The distinguishing
 feature between a shell and a plate is that the shell has a curvature whereas plates are flat structures [3]. It is, however, not necessary for the thickness
 in a shell to be extremely small. It can also vary across the surface. The behaviour of a shell can be linear or nonlinear.

Shells are ubiquitous in engineering and nature because of the advantages they present.
They carry load efficiently - this is most optimally done when all the load is carried using membrane forces; 
exhibit high structural integrity, reserved strength, strength to weight ratio; and have high stiffness [4]. 
This results in EXAMPLES HERE.

\subsection{Shell Theories: an overview}
\begin{itemize}
    \item Shell's behaviour under loading is governed primarily by its curvature
    \item Due to the curvature, their bending cannot be separated from stretching
    \item Thick shells vs thin shells.
        \begin{itemize}
            \item Thin shells are shells with a max(h/R) <= (1/20)
            \item For thick shells, (1/1000) <= h/R <= (1/20)
        \end{itemize}
    \item Shell theories aim to reduce the dimensionality of the equilibrium problem from a 3D body to only the study of the midsurface ie 2D.  
    \item Love extended the Kirchhoff hypothesis for plate bending and the set of assumptions came to be called Kirchhoff-Love Theory(?). Had its share of problems
    \item Reissner developed a linear theory for thin shells and managed to iron out some problems with the K-L theory. 
\end{itemize}

\subsection{Problems with modelling shells/membranes}
In the real world, thin shells are very sensitive to imperfections and can be highly unstable if not designed properly[5]. Modelling this complexity is challenging too. Shells can undergo large deformations, large rotations, and large strains. They can behave highly nonlinearly and
 can have multiple unstabilities - in compression, for out-of plane loading, for pressure loading, and with respect to in-plane loading [6]. 
 This mixture of performance and sensitivity makes the shell the ‘primadonna of structures’ [2]. Given that there are no analytical solutions
 for most cases to predict a shell’s behaviour, reliable approximate solutions are necessary. This report ---------.
\end{document}